%% Phd Thesis Template
%% 
%% Based  on the template of University of Wollongong
%%   2014-01-08, Thomas Griffiths, tmg994(at)uowmail.edu.au
%% Modified by 
%%   2014, Miguel Hernando Camelo
%%   2015, Ferney Maldonado, Jaime Chavarriaga

%% Preamble
%% ========

\documentclass[12pt,oneside]{book}
%% This is the `thesispreamble.tex' file,

%% Phd Thesis Template
%% ===================
%% 
%% Based  on the template of University of Wollongong
%%   2013-07-16, Thomas Griffiths, tmg994(at)uowmail.edu.au
%% Modified by 
%%   2014, Miguel Hernando Camelo
%%   2015, Ferney Maldonado, Jaime Chavarriaga

%----------------------------------------------------------------------------%
% I encourage you to read the documentation for each of the packages below,  %
% they contain instructions for implementation and examples of their use.    %
% If you get stuck read my comments, hopefully they can help you find where  %
% your answers will be. I highly reccomend making friends with Mr. Google,   %
% he knows quite a bit. The wikibook on LaTeX is also very helpful:          %
% http://en.wikibooks.org/wiki/LaTeX/                                        %
% The packages I've loaded here are the bare basics. and mainly deal with    %
% formatting, captions and things. There are thousands of packages out there %
% for all the disciplines and formatting you might need. Google what you're  %
% looking for and the keywords 'LaTeX' and 'package', You'll probably find   %
% what you're looking for. I encourage you to look on ctan, www.ctan.org,    %
% for packages that might be relevant to your degree. If you're in science I % 
% can recomend the siunitx and the chemmachro (or mhchem) package. They make %
% it really easy to typeset chemical equations, any quantity with units and  %
% scientific notation.                                                       %
%                                                                            %
%----------------------------------------------------------------------------%

\usepackage{geometry}
	\geometry{a4paper,inner=3cm, outer=3cm, top=3cm, bottom=3cm}
	% Dimensions from UOW thesis guidelines.
	\pdfpagewidth=\paperwidth 
	\pdfpageheight=\paperheight
	% This acts as a failsafe to ensure things aren't stretched or moved when it's finally printed as a PDF.

%\usepackage[parfill]{parskip} 
% Activate to begin paragraphs with an empty (return) line, comment out the indent below if you chose the return line option.

\setlength{\parindent}{4ex}	% Sets the length of the paragraph indent. Current setup has a an indent. Disable this if you activate the return line above.
	
\usepackage{setspace}
% Double or one and a half spacing.
	
\usepackage{graphicx}
	\DeclareGraphicsRule{.tif}{png}{.png}{`convert #1 `dirname #1`/`basename #1 .tif`.png}
% Graphics. Remove me and you won't have any figures, and that would be very boring.

\usepackage[usenames,dvipsnames,svgnames,table]{xcolor}
% Adds the ability to make coloured text and lines throughout the document. See documentation for xcolor.

%-------------------- Tables, figures and captions
\usepackage[font={small},labelfont={bf},margin=4ex]{caption}
% Makes bold labeled and smaller font captions. Must be loaded before the longtable package to avoid conflicts! 

\usepackage{longtable} 
% Long tables (more than one page). Different headers and footers for beginning and end pages, etc.

\usepackage{afterpage} 
% Make a longtable start on the next clear page, but fills the previous one with text first (no random gaps in the text-from long tables anymore! Man, the day I discovered this...)

\usepackage{booktabs} 
% Nice looking tables and lines in tables

\usepackage{multirow} 
% Entries in tables over multiple rows

\usepackage{lscape} 
% Pages in landscape

\usepackage{pdflscape} 
% Landscape pages also rotated in the pdf

\usepackage{wrapfig} 
% Allows figures that don't take up the entire width of the page, wrapping the text around the figure

%\usepackage[position=top,singlelinecheck=false,captionskip=4pt]{subfig} 
% Multiple figures in an individual figure. Fig. 1 a, b, c, etc. each with, or without, it's own individual caption, and with a global caption for all sub figures.

%-------------------- Special symbols and fonts
\usepackage{amssymb} 
% Maths symbols

%-------------------- Document sections, headers, footers, and bibliography

% footnotes at page bottom
\usepackage[bottom]{footmisc}
\usepackage{fancyhdr}											
% for creating different headers and footers

%-------------------- Bibliography
%\usepackage[backend=biber,articletitle=true,style=chem-rsc,doi=false]{biblatex}
%\usepackage[backend=biber,style=numeric-comp, sorting=none, clearlang=true, sortcites=true, doi=false, url=false]{biblatex}
% This is the package that lets you create a bibliography. I recommend reading the biblatex documentation to understand all the options i've specified here. BibLaTeX was created to replace BibTeX. It has lots of extra fields and options. I'm also using the biber backend here rather than the default, it copes with unicode and so can deal with accented characters easily.

% Currently this is set up to use RSC style references with article titles displayed.

% Traditionally you would use BibTeX, a special build of TeX, the newer biblatex package is a more powerful bibliograpy management tool for LaTeX. You can make multiple chapter based bibliographies, footnote bibliographes, sort your references by date, author, order cited, essentially by any bit of citation data you happen to have. You can also have a seperate library with a differnet format for say books and articles. Or if you're a PhD student, the thesis references and your publications.

%\usepackage[numbers,super,comma,sort&compress]{natbib}				
	%\setcitestyle{square}										
	% places citations in square brackets to helps to distinguish between powers and citations
%This is the old natbib package that meshes with bibtex (rather than using the newer biblatex). It's here mainly for legacy purposes. Try to shift to biblatex if you can, it is cleaner in it's implementation and creating a custom citation style is easier then with bibtex.

\usepackage[unicode=true,colorlinks=true,linkcolor=black,citecolor=black,urlcolor=black,breaklinks=true]{hyperref}
% The hyperref package allows you to have clickable links in your pdf. It also allows you to have the mailto link associated with your name. It can be  a bit finicky, so load it last.

%-------------------- Command renewals, New commands etc.
\renewcommand{\thefootnote}{\alph{footnote}}							
%letters for footnotes instead of numbers to avoid confusion with references.
 

%% Packages
%% ========

\usepackage{pdfpages}       % allows pdf inclusion
\usepackage{epstopdf}        % used to accept eps format images. 

\usepackage{amsfonts}		% includes math fonts
\usepackage{amsmath,amssymb,amsthm}
\usepackage{mathptmx}       % sets Times Roman as basic font
\usepackage{helvet}         % sets Helvetica as sans-serif font
\usepackage{courier}        % sets Courier as typewriter font
\usepackage{type1cm}        % for font management

\usepackage{makeidx}		% allows index generation
\usepackage{graphicx}		% standard LaTeX graphics tool

% two-column index, acronym lists
\usepackage{multicol, multirow, rotating,bigstrut, booktabs}
\usepackage[section]{placeins}
\usepackage[printonlyused,nohyperlinks,smaller]{acronym}

% special environments
\usepackage{algorithm, algorithmic}
\usepackage{caption}
\usepackage{subfigure}
\usepackage{relsize}
%\usepackage{subcaption}
%\usepackage[compact]{titlesec}
\newtheorem{def1}{Definition}[chapter]
\newtheorem{theo1}{Theorem}[chapter]
\newtheorem{prop1}{Proposition}[chapter]
\newtheorem{coro1}{Corollary}[chapter]
\newtheorem{lemma1}{Lemma}[chapter]

% Biblatex package
\usepackage[backend=bibtex,maxnames=4, style=numeric-comp, sorting=none, clearlang=true, sortcites=true, doi=false, url=false, firstinits]{biblatex}

\usepackage[pagewise]{lineno}
\usepackage[catalan,spanish,english]{babel}
\usepackage{csquotes}

\usepackage[PhD]{udgthesistitlepage}	% PhD style from UoW
\hyphenation{bet-ween}					% Hyphenation
\usepackage[none]{hyphenat}
\usepackage{array}						% enhanced support for tables


%% Special Settings
%% ================

\preto\chapter\acresetall		% add space after a chapter in the TOC
\def\acts{\curvearrowright}		% use "curvearrow" for group operator
 
% new column types
\newcolumntype{P}[1]{>{\centering\arraybackslash}p{#1}}
\newcolumntype{M}[1]{>{\centering\arraybackslash}m{#1}}

% Special caption for figures that includes the source
\newcommand*{\captionsource}[2]{%
  \caption[{#1}]{%
    #1%
    \\\hspace{\linewidth}%
    \textbf{Source:} #2%
  }%
}

%% BibTex files 
%% ============
%% include here your bibtext files
\addbibresource{your_bibliography.bib}
\addbibresource{own_publications.bib}
 

%% Main document configuration
%% ===========================

\begin{document}
    \sloppy				% (try to) prevent overfull boxes
    \onehalfspace		% 1.5 spacing
    %\doublespacing		% double spacing


%% Title page information
%% ======================
%% NOTE: see the uowthesistitlepage package.
    %\class{}

    \title{A PRETTY GOOD TITLE} 
    % Author: Full name, and any degrees held.
    \author{Adri\`a Piqu\'e}
	% Date: Year, alternatively use the \today macro for Month dd, yyyy.
    \date{2014} 
    % Degree: Write it in full: e.g. Bachelor of Science Medicinal Chemistry Honours
    \degree{Doctor of Philosophy by the University of Girona}
    
    % Supervisor: in square brackets: (default 1) put the number of supervisors. 
    % In the Curly braces, list your supervisor(s) seperated by commas.    
    \supervisor[2]{Ph.D. Llu\'is Vil\`a and Ph.D. Andreu Vil\`a}     
    % CoSupervisor: the same as the supervisor
    %\cosupervisor[1]{Dr. C. O. Supervisor}

    % School    
    \school{DOCTORAL PROGRAM IN TECHNOLOGY} 


%% Cover
%% =====
    %\cleardoublepage
    %\thispagestyle{empty}
{\centering
{\onehalfspacing
\newgeometry{inner=\innersidemargin,outer=\outersidemargin,top=\topsidemargin,bottom=\bottomsidemargin}
\ifsubject
    {
    \@graphic\\
    }\\ 
\else
    {
    \@graphic\\
    } 
\fi

\vspace{\biggap}
{\Large{DOCTORAL THESIS}}\\

\textbf
    {
    \vspace{\biggap}\\
    {\Large\@title}\\
    \vspace{\smallgap}
    }
\@author\\
\vspace{\biggap}

\@date\\
\vspace{\smallgap} 

\ifschool
    {
    \noindent\@school\\
    }
\else
    {
    }
\fi
\vspace{\smallgap}

\ifhorizontal
    {
    \ifcosupervisor
        {
        \begin{tabular}{cc}
        \ifmultiplesupervisors 
            {
            Supervisors:
            } 
        \else 
            {
            Supervisor: 
            }
        \fi
        &
        \ifmultiplecosupervisors 
            {
            Co-supervisors:
            } 
        \else 
            {
            Co-supervisor: 
            }\\
        \fi
        \@supervisor
        &
        \@cosupervisor
        \end{tabular}
        }\\
    \else 
        {
        \ifnosupervisor
            { 
            } 
        \else
            {
            \ifmultiplesupervisors 
                {
                Supervisors:\\ \@supervisor
                }\\ 
            \else
                {
                Supervisor:\\ \@supervisor
                }\\ 
            \fi
            } 
        \fi
        }
    \fi
    }
\else
    {
    \ifnosupervisor
        { 
        } 
    \else
        {
        \ifmultiplesupervisors 
            {
            Supervisors:\\ \@supervisor\\
            } 
        \else
            {
            Supervisor:\\ \@supervisor\\
            } 
        \fi
        } 
    \fi
    \ifcosupervisor
        {
        \ifmultiplecosupervisors 
            {
            Co-supervisors:\\ \@cosupervisor\\
            }
        \else
            {
            Co-supervisor:\\ \@cosupervisor\\
            }
        \fi
        }
    \else
        {
        } 
    \fi
    }
\fi

\vspace{\smallgap}
%\textit{
\@requirement \@degree\\
%}\\

\restoregeometry       
\clearpage{\pagestyle{empty}\cleardoublepage}
}
}


%% Front Matter
%% ============
    \makecover
    \maketitle
    %\frontmatter    
    %\linenumbers

%% Certificate
%% ===========
    \frontmatter    
    \cleardoublepage
    % Declaration
% ===========

%% To include a PDF with the scanned declaration
%% NOTE: You can define which page to include
%\includepdf[pages={1}]{figures/cert2.pdf}

%% An empty box for the demo
Here will be the declaration from the supervisors\\
\framebox(400,500){}
    \setcounter{page}{3}
    
%% Acknowledgments
%% ===============    
    \cleardoublepage
    \phantomsection
    %% Acknowledgments 
%% ===============

\chapter*{Acknowledgments}

First at all, all the credit of this thesis is to the persons that have influenced, encouraged, and helped me to the completion of this thesis. 
    \cleardoublepage\null\newpage

%% Table of contents
%% =================
    
%% List of acronyms
    \cleardoublepage
    \phantomsection
    %% List of Acronyms
%% ================

\phantomsection\addcontentsline{toc}{chapter}{List of Acronyms}
\chapter*{List of Acronyms} 

\begin{acronym}
%\setlength{\parskip}{0ex}
%\setlength{\itemsep}{0ex}
\singlespace

\acro{AuS}{Automatic Structure}
\acro{AS}{Autonomous System}
\acrodefplural{AS}[ASs]{Autonomous Systems}
\acro{WM} Windows Machine
\acro{RT}{Routing Table}
\acrodefplural{RT}[RTs]{Routing Tables}

\end{acronym}

%% List of Tables
    \cleardoublepage
    \phantomsection
    \label{listoftables}\addcontentsline{toc}{chapter}{List of Tables} 
    \listoftables
    
%% List of Figures    
    \cleardoublepage
    \phantomsection
    \label{listoffigures}\addcontentsline{toc}{chapter}{List of Figures}
    \listoffigures

%% List of Publications
    \cleardoublepage
    %% List of Publications
%% ====================

% include the references in "own_publications"
\begin{refsection}[own_publications]
    \defbibnote{myprenote}{Here we list the international journals and proceedings of international conferences where this work has been published:}
   \nocite{*}   
   \printbibliography[prenote=myprenote, heading=bibintoc, title={Related Author's Publications}]
\end{refsection}










    
%% Table of Content
    \pdfbookmark[0]{Contents}{Contents} 
    \setcounter{tocdepth}{2}
    \setcounter{secnumdepth}{3}
    \tableofcontents
    \markboth{Index}{}
    
% NOTE: These \phantomsection are to ensure that the hyperref package 
%       hyperlinks to the correct page in the electronic pdf. 
%       If you turn hyperref off they don't do anything.


%% Abstracts
%% =========

% in English
    \phantomsection	
    \cleardoublepage
    %% Abstract (in English)
%% =====================

\phantomsection\addcontentsline{toc}{chapter}{Abstract}
\chapter*{Abstract} % Starred chapter=chapter with no number.

This is the abstract for my thesis.
    
% in Spanish    
    \phantomsection	
    \cleardoublepage
    %% Abstract (in Spanish)
%% =====================

\phantomsection\addcontentsline{toc}{chapter}{Resumen}
\chapter*{Resumen} 

Este es el resumen para mi tesis.
    
% in Catalan
    \phantomsection	
    \cleardoublepage
    %% Abstract (in Catalan)
%% =====================

\phantomsection\addcontentsline{toc}{chapter}{Resum}
\chapter*{Resum} 

Aquest \'es el resum de la meva tesi.
    \mainmatter
    
%% Chapters
%% ========    
    \phantomsection	
    \cleardoublepage
    \chapter{Introduction}
\section{The Title of Your First Section}
\subsection{A Subsection --- for Clarity}
Lorem ipsum dolor sit amet, consectetur adipiscing elit.\cite{pericles} Integer ut congue lectus. Nullam dapibus scelerisque diam, ac convallis dolor convallis non. Sed nec lectus nec sapien interdum commodo nec quis elit. Sed neque augue, pulvinar id imperdiet id, ullamcorper sed sem. Duis pulvinar blandit erat, quis suscipit erat venenatis sit amet. Aenean varius aliquam dignissim. Sed tempus consequat sapien et bibendum. Class aptent taciti sociosqu ad litora torquent per conubia nostra, per inceptos himenaeos. Maecenas arcu ligula, tincidunt non volutpat nec, luctus ac justo. Donec euismod egestas leo, vel convallis mi accumsan sed. Aliquam erat volutpat. Maecenas volutpat lacinia justo, dictum iaculis enim scelerisque in. Vestibulum consequat augue in nisl luctus eget ultrices sem blandit. Donec mi risus, rutrum at tempus in, dapibus vel est. Sed eu ullamcorper velit.\cite{pericles,thekingsenglish}

Duis malesuada ultrices rutrum. In posuere sem dapibus urna accumsan id sollicitudin turpis iaculis. Sed fermentum metus vel dui vulputate consectetur. Nunc nec ante nisi, ac ultricies metus. Nullam sed dui vitae metus vehicula adipiscing. Nunc facilisis tortor neque. Praesent ac nulla at odio lacinia egestas. Nam facilisis vehicula pretium.

Fusce tempus libero vitae leo cursus cursus. Suspendisse iaculis dignissim placerat. Vivamus ornare, lectus in placerat cursus, quam ligula congue risus, a facilisis leo est ut massa. Nulla in metus id eros rutrum aliquam. Cras id velit ut enim iaculis adipiscing facilisis sit amet lorem. Ut id neque id turpis facilisis tempus sit amet nec urna.\cite{bioluminescence:ch1} Ut volutpat magna et dolor elementum eget elementum lorem eleifend. Pellentesque mollis convallis est a tincidunt. Nunc egestas, ante a blandit iaculis, risus orci auctor risus, in pellentesque est sapien id mi. Nunc faucibus porttitor tincidunt.\cite{atkins_inorgchem,atkins_physchem} Fusce eu faucibus dolor. Mauris quis arcu metus, at malesuada lacus. Nam lectus lacus, sodales nec tempus eu, pretium id risus. Cras dignissim aliquet laoreet. Morbi sed arcu id tellus tempor consectetur.\cite{bioluminescence:apA,hori1973}

\subsection{Another Subsection}
Aenean convallis, ante quis convallis viverra, urna lorem venenatis odio, ut bibendum nisi mauris a mi. Nam laoreet arcu a eros dignissim elementum. Class aptent taciti sociosqu ad litora torquent per conubia nostra, per inceptos himenaeos. Proin vitae magna ligula, id consequat felis. Integer rutrum ante vitae tellus convallis volutpat in vitae odio. Vivamus eu turpis et turpis euismod faucibus. Suspendisse at turpis sit amet magna auctor fringilla. Vestibulum quis nulla sed dui pellentesque dapibus.\footnote{This is a footnote. Footnotes are notes at the foot of the page. Literary style guides (and some supervisors) recommend limited use of footnotes; but, publishers often encourage them. Use them as you will but don't go overboard.} Cum sociis natoque penatibus et magnis dis parturient montes, nascetur ridiculus mus. Maecenas id odio nisl. Etiam facilisis elit in urna hendrerit ac iaculis mauris lacinia. Fusce congue ultricies nulla vel elementum. Phasellus suscipit vestibulum tortor at semper. Etiam eu sagittis augue.\cite{bioluminescence,clayden_orgchem}


\begin{figure}
\centering
\includegraphics[height=0.45\textwidth]{figures/example.jpg}
\caption[Chemistry cat.]{This is chemistry cat. Here he serves to demonstrate a figure in a LaTeX document, complete with caption. Notice his suave bow tie, but also that he has forgotten to label his solutions.}
\end{figure}

Duis malesuada ultrices rutrum. In posuere sem dapibus urna accumsan id sollicitudin turpis iaculis. Sed fermentum metus vel dui vulputate consectetur. Nunc nec ante nisi, ac ultricies metus. Nullam sed dui vitae metus vehicula adipiscing. Nunc facilisis tortor neque. Praesent ac nulla at odio lacinia egestas. Nam facilisis vehicula pretium.

\begin{figure}[thb]
	\centering
		\includegraphics[width=0.45\textwidth]{figures/example.jpg}
    	\captionsource{This is the same \ac{RT} cat. Note the additional field for the source.}{http://bgp.potaroo.net.}
	\label{fig:fig2}
\end{figure}

Praesent sed diam arcu, quis aliquam libero \ac{AuS}. Pellentesque a nulla vulputate lectus rhoncus vulputate ut in urna \ac{AS}. Quisque sapien sem, convallis ac porta sit amet, aliquet sit amet eros. Etiam quis lorem ligula. Mauris elit sapien, ultrices nec vulputate ut, porta eu leo. Ut consequat accumsan commodo. Sed ut ultricies ante.

\begin{figure}[thb]
	\centering
		\includegraphics[width=0.45\textwidth]{figures/example.jpg}
    	\captionsource[The same cat]{This is the same \ac{RT} cat. Note the
    	additional field for the source.}{http://bgp.potaroo.net.}
	\label{fig:fig3}
\end{figure}

Duis malesuada ultrices rutrum. In posuere sem dapibus urna accumsan id
sollicitudin turpis iaculis. Sed fermentum metus vel dui vulputate consectetur. Nunc nec ante nisi, ac ultricies metus. Nullam sed dui vitae metus vehicula adipiscing. Nunc facilisis tortor neque. Praesent ac nulla at odio lacinia egestas. Nam facilisis vehicula pretium.

\begin{algorithm}\label{Alg1}
	\begin{algorithmic}[1] 
	 \caption{Greedy embedding in \ac{WM} spaces.}
	\label{alg:AlgorithmX}
				\STATE Choose an arbitrary vertex $r\in{V(H)}$.\\
				%\STATE Use the \ac{BFS} algorithm to compute a spanning tree $T_{H}$ rooted in $r$.\\
                \STATE Use a distributed algorithm to compute a spanning tree $T_{H}$ rooted at $r$.
				\STATE Compute the maximum degree $\Delta_H$ of the graph.\\
				\STATE For each $v\in{V(H)}$, enumerate its children and assign them a unique integer $i\in{\{1, \dots{}, \Delta_H\}}$.\\
				\STATE Assign to each vertex a unique word (label) representing a group element from a specific algebraic group $G$.\\ 
	\end{algorithmic}
\end{algorithm}


Praesent sed diam arcu, quis aliquam libero \ac{AuS}. Pellentesque a nulla vulputate lectus rhoncus vulputate ut in urna \ac{AS}. Quisque sapien sem, convallis ac porta sit amet, aliquet sit amet eros. Etiam quis lorem ligula. Mauris elit sapien, ultrices nec vulputate ut, porta eu leo. Ut consequat accumsan commodo. Sed ut ultricies ante.

\begin{table}[H]\scriptsize
  \centering
    \begin{tabular}{|M{4cm}|M{2.5cm}|M{2.5cm}|M{2cm}|M{2cm}|}
    \hline
    \multirow{2}[4]{*}{\textbf{Spanning Tree}} & \multirow{2}[4]{*}{\textbf{Time}} & \multirow{2}[4]{*}{\textbf{Space}} & \multicolumn{2}{c|}{\textbf{Message}} \bigstrut\\
\cline{4-5}                &       &       & \textbf{Size} & \textbf{Number} \bigstrut\\
    \hline
    \ac{AuS} & $O(D+log^2(n))$ & $O(\Delta_H\cdot{}log(n))$ & $O(log(n))$ & $O(m\cdot{}log^2(n))$ \bigstrut\\
    \hline
    \ac{AuS}  & $O(D_H)$  & $O(\Delta_H)$ & $O(1)$  & $O(D\cdot{}m)$ \bigstrut\\
    \hline
    \ac{AS}   & $O(n)$  & $O(\Delta_H\cdot{}log(n))$ & $O(n)$  & $O(n)$ \bigstrut\\
    \hline
    \end{tabular}%
  \caption{Strategies for distributed computing of a rooted spanning tree from an undirected, unweigthed and connected graph.}
  \label{tab:spanning-trees}%
\end{table}%

Duis malesuada ultrices rutrum. In posuere sem dapibus urna accumsan id sollicitudin turpis iaculis. Sed fermentum metus vel dui vulputate consectetur. Nunc nec ante nisi, ac ultricies metus. Nullam sed dui vitae metus vehicula adipiscing. Nunc facilisis tortor neque. Praesent ac nulla at odio lacinia egestas. Nam facilisis vehicula pretium.

    \chapter{Your Second Chapter}

    \chapter{Your Third Chapter}

    \chapter{Your Fourth Chapter}

    \chapter{Your Fifth Chapter}


% NOTE: The chapters can be included using \input or \include
% \include creates an aux file for the included file, \input does not. 
% \include is good for large files, like a chapter, since LaTeX won't re-compile 
% the file if no changes have been made. \input is better for small inputs like 
% data tables. Additionally you *CAN* nest \input commands. 
% You *CANNOT* nest \include commands.


%% Bibliography
%% ============
    \cleardoublepage
    \phantomsection	
    \addcontentsline{toc}{chapter}{Bibliography}    
    \printbibliography

%% Appendices
%% ==========
    \cleardoublepage
    \appendix
    \chapter{Title for the Appendix One}
\section{The title of the first section}

The appendices work exactly the same way as chapters, they are numbered with letters rather than numbers though.
    \chapter{Title for the Appendix Two}
\section{The title of the first section}
    
\end{document}  

%% ----------------------------------------------------------

%% Notes and Examples
%% ==================

%% Figure or Scheme (caption after)

% \begin{figure}
% \centering
% \input{figures/name_of_figure.tex}
% \caption[Short caption for table of figures]{Long caption for text body}
% \end{figure}

%% Table (caption before)

% \begin{table}
% \centering
% \caption[Short caption for table of tables]{Long caption for text body}
% \input{figures/name_of_figure.tex}
% \end{table}

